\subsection{Structural Abstraction of the Unicode Standard}

A \texttt{str} is a sequence of Unicode codepoints, not a byte buffer and not a rendered glyph. \texttt{str.encode()} crosses to \texttt{bytes}; \texttt{bytes.decode()} returns to text. The two types expose different method sets (\texttt{encode} vs.\ \texttt{decode}, \texttt{hex} on bytes only).

\subsubsection{Distinction Between Abstract Codepoints, Visible Glyphs, and Transmitted Byte Serializations}

\texttt{ord()} reads the codepoint number of a one-character \texttt{str}. Encoding with UTF-8 produces a variable-length byte sequence per codepoint; \texttt{len(ch)} counts codepoints, not UTF-8 width. Glyphs can differ from codepoint count: \texttt{"é"} as \texttt{U+00E9} versus \texttt{"e\textbackslash u0301"} (base letter plus combining acute) may render similarly but compare unequal and normalize differently under NFC. Python indexes and stores the codepoint sequence; rendering and byte serialization are downstream concerns.

\subsubsection{Text vs. Bytes: Why \texttt{str} Stores Text and \texttt{bytes} Stores Raw Byte Values}

Indexing \texttt{str} yields a one-character \texttt{str}; indexing \texttt{bytes} yields an \texttt{int} in \texttt{0--255}. Iteration mirrors this split. Python refuses \texttt{str + bytes} because no default encoding is implied---conversion must be explicit. At I/O boundaries, programs decode incoming \texttt{bytes} to \texttt{str} for text semantics and encode outgoing \texttt{str} when the protocol requires octets.
